% GENERATED FILE. Edit canonical lessons, then run npm run generate:books.
% canonical-source-hash: fnv1a64:646bdc590b3569f5
% canonical-lessons: RU-C216-spokoynyy, RU-C216-nervnyy, RU-C216-schastlivyy, RU-C216-grustnyy, RU-C216-strogiy

\chapter{Calm, Nervous, Happy, Sad, Strict}
\label{ch:ru-calm-nervous-happy-sad-strict}
\hlchaptermodality{216}

\begin{chapteropening}
\textbf{By the end of this chapter:} \emph{I can say calm, nervous, happy, sad and strict.}

Complete the last lesson of chapter 216: I can say calm, nervous, happy, sad and strict.
\end{chapteropening}

\section[\texorpdfstring{\ru{спокойный} (spokóynyy)}{спокойный (spokóynyy)}]{\ru{спокойный} (spokóynyy) — calm}

\label{lesson:RU-C216-spokoynyy}



\begin{quote}
Before the new one: say the Russian for unfortunately, then the Russian for fortunately.
\end{quote}

\subsection*{You'll want to know: \ru{спокойный}}
\textbf{\ru{спокойный}} (\emph{spokóynyy}) — \textquotedblleft{}calm\textquotedblright{}.

\begin{cousinweb}
\textbf{\ru{спокойный}} is built on \emph{\ru{покой}}, \textquotedblleft{}rest, peace\textquotedblright{}, which Vasmer relates to Latin \emph{quiēs}, \textquotedblleft{}rest\textquotedblright{}, and Gothic \emph{hveila}, \textquotedblleft{}a while\textquotedblright{}: English \emph{while} is a cousin.
\end{cousinweb}

\begin{grammarlens}[title={What you've built}]
One more word for this chapter.
\end{grammarlens}

\subsection*{Your turn}
\begin{itemize}
\raggedright
  \item \emph{Say it:} \emph{spokóynyy}
  \item \emph{Say it:} \emph{spokóynyy}, once more
  \item \emph{Say it:} say \emph{spokóynyy}
  \item \emph{Recall:} say the Russian for unfortunately, then the Russian for fortunately, then say \emph{spokóynyy} again
\end{itemize}

\begin{tcolorbox}[breakable,colback=teal!4,colframe=teal!35!black,title={Before you move on}]
What does \ru{спокойный} mean? (\textquotedblleft{}Calm\textquotedblright{}.) Say it once more.
\end{tcolorbox}
\section[\texorpdfstring{\ru{нервный} (nérvnyy)}{нервный (nérvnyy)}]{\ru{нервный} (nérvnyy) — nervous}

\label{lesson:RU-C216-nervnyy}



\begin{quote}
Before the new one: say the Russian for fortunately, then the Russian for calm.
\end{quote}

\subsection*{You'll want to know: \ru{нервный}}
\textbf{\ru{нервный}} (\emph{nérvnyy}) — \textquotedblleft{}nervous\textquotedblright{}.

\begin{cousinweb}
\textbf{\ru{нервный}} is \emph{\ru{нерв}}, \textquotedblleft{}a nerve\textquotedblright{}, with \emph{-\ru{ный}}. \emph{\ru{нерв}} is a bookish loan through German \emph{Nerv} or French \emph{nerf} from Latin \emph{nervus}, \textquotedblleft{}a sinew, a nerve\textquotedblright{} (Vasmer).
\end{cousinweb}

\begin{grammarlens}[title={What you've built}]
One more word for this chapter.
\end{grammarlens}

\subsection*{Your turn}
\begin{itemize}
\raggedright
  \item \emph{Say it:} \emph{nérvnyy}
  \item \emph{Say it:} \emph{nérvnyy}, once more
  \item \emph{Say it:} say \emph{nérvnyy}
  \item \emph{Recall:} say the Russian for fortunately, then the Russian for calm, then say \emph{nérvnyy} again
\end{itemize}

\begin{tcolorbox}[breakable,colback=teal!4,colframe=teal!35!black,title={Before you move on}]
What does \ru{нервный} mean? (\textquotedblleft{}Nervous\textquotedblright{}.) Say it once more.
\end{tcolorbox}
\section[\texorpdfstring{\ru{счастливый} (schastlívyy)}{счастливый (schastlívyy)}]{\ru{счастливый} (schastlívyy) — happy}

\label{lesson:RU-C216-schastlivyy}



\begin{quote}
Before the new one: say the Russian for calm, then the Russian for nervous.
\end{quote}

\subsection*{You'll want to know: \ru{счастливый}}
\textbf{\ru{счастливый}} (\emph{schastlívyy}) — \textquotedblleft{}happy\textquotedblright{}.

\begin{cousinweb}
\textbf{\ru{счастливый}} is built on \emph{\ru{счастье}}, \textquotedblleft{}luck, happiness\textquotedblright{}, usually explained as \textquotedblleft{}a good share\textquotedblright{}, from \emph{\ru{часть}}, \textquotedblleft{}a share\textquotedblright{} (Vasmer). Vasmer relates \emph{\ru{часть}} to Lithuanian \emph{kąsti}, \textquotedblleft{}to bite\textquotedblright{}: a share is a bite.
\end{cousinweb}

\begin{grammarlens}[title={What you've built}]
One more word for this chapter.
\end{grammarlens}

\subsection*{Your turn}
\begin{itemize}
\raggedright
  \item \emph{Say it:} \emph{schastlívyy}
  \item \emph{Say it:} \emph{schastlívyy}, once more
  \item \emph{Say it:} say \emph{schastlívyy}
  \item \emph{Recall:} say the Russian for calm, then the Russian for nervous, then say \emph{schastlívyy} again
\end{itemize}

\begin{tcolorbox}[breakable,colback=teal!4,colframe=teal!35!black,title={Before you move on}]
What does \ru{счастливый} mean? (\textquotedblleft{}Happy\textquotedblright{}.) Say it once more.
\end{tcolorbox}
\section[\texorpdfstring{\ru{грустный} (grústnyy)}{грустный (grústnyy)}]{\ru{грустный} (grústnyy) — sad}

\label{lesson:RU-C216-grustnyy}



\begin{quote}
Before the new one: say the Russian for nervous, then the Russian for happy.
\end{quote}

\subsection*{You'll want to know: \ru{грустный}}
\textbf{\ru{грустный}} (\emph{grústnyy}) — \textquotedblleft{}sad\textquotedblright{}.

\begin{cousinweb}
\textbf{\ru{грустный}} is \emph{\ru{грусть}}, \textquotedblleft{}sadness\textquotedblright{}, with \emph{-\ru{ный}}. Vasmer reports the usual link of \emph{\ru{грусть}} with Lithuanian \emph{graudùs}, \textquotedblleft{}brittle; sad\textquotedblright{}.
\end{cousinweb}

\begin{grammarlens}[title={What you've built}]
One more word for this chapter.
\end{grammarlens}

\subsection*{Your turn}
\begin{itemize}
\raggedright
  \item \emph{Say it:} \emph{grústnyy}
  \item \emph{Say it:} \emph{grústnyy}, once more
  \item \emph{Say it:} say \emph{grústnyy}
  \item \emph{Recall:} say the Russian for nervous, then the Russian for happy, then say \emph{grústnyy} again
\end{itemize}

\begin{tcolorbox}[breakable,colback=teal!4,colframe=teal!35!black,title={Before you move on}]
What does \ru{грустный} mean? (\textquotedblleft{}Sad\textquotedblright{}.) Say it once more.
\end{tcolorbox}
\section[\texorpdfstring{\ru{строгий} (strógiy)}{строгий (strógiy)}]{\ru{строгий} (strógiy) — strict}

\label{lesson:RU-C216-strogiy}



\begin{quote}
Before the new one: say the Russian for happy, then the Russian for sad.
\end{quote}

\subsection*{You'll want to know: \ru{строгий}}
\textbf{\ru{строгий}} (\emph{strógiy}) — \textquotedblleft{}strict\textquotedblright{}.

\begin{cousinweb}
\textbf{\ru{строгий}} — it is usually taken as a loan from Polish \emph{srogi}, \textquotedblleft{}severe\textquotedblright{}, and the whole group is linked with \emph{\ru{стеречь}}, \textquotedblleft{}to guard\textquotedblright{}, and \emph{\ru{сторож}}, \textquotedblleft{}a watchman\textquotedblright{} (Vasmer); Trubachev's note cites other views.
\end{cousinweb}

\begin{grammarlens}[title={What you've built}]
One more word for this chapter.
\end{grammarlens}

\subsection*{Your turn}
\begin{itemize}
\raggedright
  \item \emph{Say it:} \emph{strógiy}
  \item \emph{Say it:} \emph{strógiy}, once more
  \item \emph{Say it:} say \emph{strógiy}
  \item \emph{Recall:} say the Russian for happy, then the Russian for sad, then say \emph{strógiy} again
\end{itemize}

\begin{tcolorbox}[breakable,colback=teal!4,colframe=teal!35!black,title={Before you move on}]
What does \ru{строгий} mean? (\textquotedblleft{}Strict\textquotedblright{}.) Say it once more.
\end{tcolorbox}
